% Arbeitsjournal -- separates Abgabedokument.
% Build: tectonic arbeitsjournal.tex
\documentclass[11pt,a4paper]{scrartcl}

\usepackage[T1]{fontenc}
\usepackage{lmodern}
\usepackage[nswissgerman]{babel}
\usepackage{csquotes}
\usepackage{array}
\usepackage{booktabs}
\usepackage{longtable}
\usepackage[margin=2.2cm]{geometry}

\begin{document}

\begin{center}
  {\LARGE\bfseries Arbeitsjournal}\\[4pt]
  {\large Transferarbeit - Raphael Bolliger}\\[2pt]
  CAS Machine Learning, Hochschule Luzern -- Informatik
\end{center}

\vspace{0.5em}

\begin{longtable}{@{}l r >{\raggedright\arraybackslash}p{\dimexpr\textwidth-6.4cm\relax}@{}}
\toprule
\textbf{Datum} & \textbf{Dauer (h)} & \textbf{Tätigkeit} \\
\midrule
\endfirsthead
\toprule
\textbf{Datum} & \textbf{Dauer (h)} & \textbf{Tätigkeit} \\
\midrule
\endhead

Sa, 25.07.2026 & 5.0 &
Recherche zu Objekterkennung und YOLO, Versuche mit Roboflow, Entscheid für
synthetische Trainingsdaten. \\
So, 26.07.2026 & 7.0 &
Französisches Kartenblatt fotografiert, Erkennungsziel festgelegt und
Grundlagen des Datengenerators umgesetzt. \\
Mo, 27.07.2026 & 8.0 &
Datengenerator und Labelvarianten A, B und C erweitert, Dataset Tool begonnen
und reale Aufnahmen gelabelt. \\
Di, 28.07.2026 & 8.0 &
Datensatzstruktur, Trainingsskripte und CI umgesetzt, erste Modelle trainiert
und iOS-Prototyp getestet. \\
Mi, 29.07.2026 & 8.0 &
Datengenerator erweitert, reale Negativbilder ergänzt, Cloud-Pipeline
aufgebaut und weitere Modelle trainiert. \\
Do, 30.07.2026 & 8.0 &
Scans und Labels korrigiert, flache Aufnahmen gelabelt sowie
Variantenvergleich und Modellexport weitergeführt. \\
Fr, 31.07.2026 & 6.0 &
Variantenvergleich abgeschlossen, vollständige Verarbeitung von B geprüft
und Konzept der Zählsession erstellt. \\
So, 02.08.2026 & 3.0 &
Kernablauf der Zählsession umgesetzt und auf dem iPhone geprüft. \\
Sa, 08.08.2026 & 5.0 &
Erkennung unter realen Bedingungen getestet, Fehlerfälle aufgenommen und
Stabilitätsregel geprüft. \\
So, 09.08.2026 & 2.0 &
Alternative Stabilitätsregel umgesetzt und Aufzeichnung einer Zählsession
vorbereitet. \\
So, 06.09.2026 & 8.0 &
Cloud-Pipeline überarbeitet, Winkelverteilung festgelegt, Azure eingerichtet
und Trainingsläufe gestartet. \\
Mi, 09.09.2026 & 5.0 &
Beide Kartenblätter in Generator und App ergänzt, Trainingsläufe geprüft. \\
Do, 10.09.2026 & 3.0 &
Training mit 200'000 Bildern ausgewertet, Farbfehler und Fehlalarme untersucht. \\
Fr, 11.09.2026 & 3.0 & Dokumentation erweitert. \\
Sa, 12.09.2026 & 6.0 &
Deutsches Kartenblatt aufgenommen und gelabelt, Dataset Tool erweitert und
Android-Prototyp getestet. \\
So, 13.09.2026 & 4.0 &
Android-Umsetzung finalisiert, finaler Trainingslauf durchgeführt. \\
Mo, 14.09.2026 & 6.0 &
Session-Analyse umgesetzt, Konfidenzfehler korrigiert, Schlussmodelle
ausgewertet und Dokumentation. \\
So, 20.09.2026 & 6.0 & Finalisieren der Arbeit. \\
\midrule
\textbf{Total} & \textbf{101.0} &  \\
\bottomrule
\end{longtable}

\vspace{2em}
\noindent
Ort, Datum: \rule{5cm}{0.4pt} \hfill Unterschrift: \rule{5cm}{0.4pt}

\end{document}
